\chapter{Introduction} \label{Introduction}

\section{Background and Motivation}

\label{Background and Motivation}

Large Language Models (LLMs) have fundamentally reshaped the landscape of artificial intelligence, demonstrating remarkable capabilities across natural language understanding, generation, and reasoning tasks. Building on this momentum, the field has increasingly moved toward Multimodal Large Language Models (MLLMs) that extend these capabilities beyond text to visual domains, enabling models to understand, generate, and manipulate images alongside language \cite{liu2023visualinstructiontuning, openai2024gpt4technicalreport}. Among the most rapidly advancing areas is AI-driven image generation and editing, which has seen dramatic improvements in both quality and controllability over recent years \cite{rombach2022highresolutionimagesynthesislatent, zhang2023addingconditionalcontroltexttoimage}.

A particularly consequential development has been the emergence of text-guided image editing, where users can modify images through natural language instructions rather than through manual, tool-based workflows. This paradigm shift lowers the barrier to image manipulation: tasks that previously required professional expertise in tools like Adobe Photoshop can now be performed by specifying the desired change in words \cite{adobe2023firefly}. However, not all image editing tasks are equal in difficulty. A useful distinction can be drawn between two categories of image editing that differ in how success is measured. We define qualitative editing as modifications to the semantic content of an image where correctness is assessed subjectively or categorically, for example, changing an object's style, replacing a background, or altering the mood of a scene. In these tasks, multiple outputs may be considered equally valid, and evaluation typically relies on perceptual quality or human preference. By contrast, we define quantitative editing as transformations that involve precise, numerically specified changes where the target output is uniquely determined, for example, repositioning a text element by exactly 40 pixels, scaling a heading to 150\% of its original size, or changing a color to a specific hex value. In quantitative editing, success can be measured against a single ground truth with pixel-level or numerical precision. While qualitative image editing has seen substantial progress and is approaching practical usability in commercial products \cite{adobe2023firefly}, quantitative editing remains a significantly harder and less explored problem.

Within image editing, text content manipulation, that is, changing what text says in an image, has received considerable attention and has seen meaningful progress in recent years. Recent work has demonstrated that text rendering and content replacement in images are increasingly reliable, with both academic models \cite{tuo2024anytext} and commercial systems such as Nano Banana Pro \cite{google2025nanobanana} showing strong performance on these tasks. However, quantitative text editing in images, such as scaling text by a precise amount or repositioning text elements within a layout remains a largely unsolved challenge. The TextEditBench benchmark \cite{gui2025texteditbench} has highlighted that current state-of-the-art models still struggle significantly with these quantitative aspects of text editing, even as they excel at content-level changes.

Understanding why these tasks remain difficult requires examining the fundamental limitations of current model architectures. Diffusion-based models, which form the backbone of many leading image generation and editing systems, suffer from attention leakage during the denoising process, which has been shown to be particularly problematic for text regions, making it difficult for these models to perform spatially precise, layout-aware edits without unintentionally modifying surrounding regions \cite{shi2025wordcon}. Meanwhile, the Visual Autoregressive (VAR) architecture, which employs next-scale prediction to generate images at progressively higher resolutions, faces its own set of challenges. Because VAR models build images from coarse to fine scales, small quantitative changes such as a two-pixel shift may not be clearly representable or distinguishable at the lower resolution scales where global structure is determined \cite{tian2024var}. These architectural constraints suggest that different model families may exhibit fundamentally different failure modes when confronted with quantitative editing tasks, a distinction that becomes central to our comparative analysis.

Compounding these architectural challenges is a severe data scarcity problem. Training data for quantitative image editing, specifically, paired examples of before-and-after images with precise, measurable transformations, is difficult to collect at scale from natural sources. Existing datasets for image editing tend to focus on qualitative changes and lack the pixel-perfect ground truth needed to train and evaluate models on exact spatial and numerical transformations. At the same time, it is well established that synthetic data can be a powerful tool for improving MLLM performance in data-scarce domains, and that even relatively modest amounts of high-quality training data can lead to state-of-the-art results in image editing \cite{zhang2025icedit}. Can carefully constructed synthetic data help bridge the gap in quantitative text editing performance?

In this work, we propose to address the quantitative text editing problem through a synthetic data pipeline that leverages browser rendering of programmatically generated HTML and CSS to produce pixel-perfect ground truth image pairs. By using an LLM to systematically generate HTML/CSS code that represents text editing operations, such as resizing, repositioning, recoloring, and content swapping, and then rendering these through a browser engine, we can produce training data with exact, verifiable ground truth at scale. 

This approach allows us not only to train models but also to construct fine-grained evaluations that isolate specific failure modes, including spatial relocation accuracy, quantitative scaling precision, color change fidelity, text content swap correctness, and unintended region modifications. We apply this data to fine-tune two models based on different architectures, a diffusion-based model and a VAR-based model, with the goal of both improving overall quantitative text editing performance and characterizing how architectural differences manifest as distinct failure mode patterns. We further evaluate the resulting models on the established TextEditBench benchmark \cite{gui2025texteditbench} to assess whether the synthetic training data yields improvements that generalize beyond our controlled evaluation setting.

\section{Research Questions}

\label{Research Questions}

\textbf{RQ1:} Can synthetic data with pixel-perfect ground truth improve MLLM performance on quantitative text editing tasks?

\textbf{RQ1a:} Do these improvements from synthetic data training generalize to established third-party benchmarks?

\noindent \textbf{RQ2:} Do diffusion-based and autoregressive (VAR) architectures exhibit distinct, architecture-specific failure modes on quantitative text editing when trained on the same data?

\section{Scope and Boundaries}

\label{Scope and Boundaries}

\textit{Script and language.} We restrict our scope to Latin-script text. While quantitative editing challenges such as repositioning and scaling are in principle script-agnostic, evaluation tooling, particularly OCR-based validation, is substantially more reliable and accessible for Latin characters. Extending the pipeline and evaluation to logographic scripts (e.g., Chinese, Japanese) or right-to-left scripts (e.g., Arabic, Hebrew) would introduce confounding variables related to text rendering and recognition accuracy that are outside the focus of this work.

\textit{Image domain.} We target digitally designed artifacts such as posters, infographics and presentations rather than natural photographs with incidental scene text (e.g., street signs, storefronts). This choice reflects both the primary practical use case for quantitative text editing and the visual domain of our synthetic training data, which is browser-rendered HTML/CSS. Generalization to natural scene text is not a goal of this thesis.

\textit{Edit types.} We focus on four categories of quantitative text editing: spatial repositioning, scaling (resizing), color modification, and content replacement with preserved layout. We do not address composite edits that combine multiple operations in a single instruction (e.g., "make the title red and move it to the top"), nor do we target edits that require semantic reasoning about image content (e.g., "make the subtitle match the tone of the image").

\textit{Evaluation scope.} Our diagnostic evaluation measures editing precision on isolated, atomic operations. We do not evaluate broader editing capabilities such as multi-turn editing workflows, undo/redo consistency, or interactive refinement. The external benchmark evaluation on TextEditBench follows that benchmark's standard protocol and is included for generalizability assessment, not as a primary success criterion.

\textit{Models.} We compare at most one model per architectural family (diffusion and VAR). Our findings about architecture-specific failure modes are therefore specific to the chosen model instances and should be interpreted as indicative rather than exhaustive characterizations of their respective architectural families.

